\documentclass[11pt]{article}
\usepackage[margin=1in]{geometry}
\usepackage{amsmath,amssymb}
\usepackage{xcolor}
\usepackage{microtype}
\usepackage{hyperref}
\setlength{\emergencystretch}{2em}

\hypersetup{colorlinks=true,linkcolor=blue!55!black,citecolor=blue!55!black,urlcolor=blue!55!black}

\title{A Negative Correction to the Proposed Inverse-Series Pattern}
\author{Literature-implied status note for arXiv:2010.00746}
\date{21 July 2026}

\begin{document}
\maketitle

\begin{center}
\fcolorbox{blue!55!black}{blue!4}{\parbox{0.90\textwidth}{
\textbf{Status.} Literature-implied negative correction.  Research Problem
1's proposed common-coefficient representation is false already at $n=5$.
Oertel's expanded sequel arXiv:2305.04428 silently replaces it by the correct
multinomial Bell-polynomial and determinant formulas.}}
\end{center}

\section*{Original problem}

Section 5.1 of Oertel's 2020 paper, source PDF pages 13--15, first recalls the
Lagrange-inversion formula for the coefficients of
\[
 h(x)=\sum_{j\geq1}\alpha_jx^j,\qquad
 h^{-1}(y)=\sum_{n\geq1}\beta_ny^n,
 \quad \alpha_1\neq0.
\]
It then proposes that the middle blocks can be written in the form
\begin{equation}\label{eq:proposal}
 n\beta_n\alpha_1^{2n-1}
 =-n\alpha_1^{n-2}\alpha_n
 +\sum_{k=2}^{n-2}(-1)^k m_k
 \binom{n-1+k}{k}\alpha_1^{n-1-k}p_k
 +(-1)^{n-1}\binom{2n-2}{n-1}\alpha_2^{n-1},
\end{equation}
where $m_k\in\{1,\ldots,k\}$ and $p_k$ is a sum of distinct admissible
monomials, each with coefficient one.  Item (1-RP1) asks to prove this
representation and determine the sets of monomials and the numbers $m_k$
\cite{oertel2020}.

\section*{The $n=5$ contradiction}

The same source explicitly lists
\begin{align*}
5\beta_5\alpha_1^9={}&-5\alpha_1^3\alpha_5
+30\alpha_1^2\alpha_2\alpha_4
+15\alpha_1^2\alpha_3^2\\
&-105\alpha_1\alpha_2^2\alpha_3+70\alpha_2^4.
\end{align*}
For $n=5$ and $k=2$, formula \eqref{eq:proposal} supplies the single block
\[
 m_2\binom{6}{2}\alpha_1^2p_2=15m_2\alpha_1^2p_2.
\]
Both $\alpha_2\alpha_4$ and $\alpha_3^2$ must occur in $p_2$, since no other
block has the same power of $\alpha_1$ and the same weighted degree.  Because
$p_2$ assigns coefficient one to every selected monomial, the proposed form
assigns the same coefficient $15m_2$ to both.  The actual coefficients are
$30$ and $15$.  No single $m_2$ can equal both $2$ and $1$, so (1-RP1) is
false as printed.

\section*{Correct formula and later source}

The correct formula retains the multinomial weights:
\begin{equation}\label{eq:correct}
 n\beta_n\alpha_1^{2n-1}
 =\sum_{k=1}^{n-1}(-1)^k\binom{n-1+k}{k}
 \alpha_1^{n-1-k}B^{\circ}_{n-1,k}
 (\alpha_2,\ldots,\alpha_{n-k+1}).
\end{equation}
Indeed,
\[
B^{\circ}_{4,2}(\alpha_2,\alpha_3,\alpha_4)
=2\alpha_2\alpha_4+\alpha_3^2,
\]
because the multiplicities of the partitions $4=1+3$ and $4=2+2$ are
$2!/(1!1!)=2$ and $2!/2!=1$.  Multiplication by
$\binom62=15$ gives exactly $30$ and $15$.

Section 9.1 of Oertel's expanded sequel, supporting PDF pages 116--120,
states \eqref{eq:correct}, emphasizes the multinomial coefficients in
$B^{\circ}_{n-1,k}$, and supplies recursive and determinant representations
\cite{oertel2023}.  It no longer contains the uniform-$m_k$ proposal.  The
later paper does not explicitly say that (1-RP1) was false; the negative
answer is therefore a direct identification from its corrected formula.

\section*{Scope and search record}

This refutes only the proposed simplification in (1-RP1).  It does not affect
the valid Lagrange--Buermann formula or the broader questions concerning
computable inverse coefficients and Grothendieck bounds.  The run indexes and
a bounded web search on 21 July 2026 found no exact published discussion of
the failed common-coefficient ansatz.  The later paper was decisive for
classifying the observation as an implication of the later literature,
rather than as a new counterexample.

\begin{thebibliography}{9}

\bibitem{oertel2020}
F. Oertel,
\emph{Grothendieck's inequality and completely correlation preserving
functions -- a summary of recent results and an indication of related
research problems}, arXiv:2010.00746v2, 2020.

\bibitem{oertel2023}
F. Oertel,
\emph{Upper bounds for Grothendieck constants, quantum correlation matrices
and CCP functions}, arXiv:2305.04428v4, revised 2025.

\end{thebibliography}

\end{document}
